\documentclass[11pt]{article}
\usepackage[T1]{fontenc}
\usepackage[utf8]{inputenc}
\usepackage{lmodern}
\usepackage[margin=1.15in]{geometry}
\usepackage{microtype}
\usepackage{setspace}
\PassOptionsToPackage{hyphens}{url}
\usepackage[hidelinks]{hyperref}
\Urlmuskip=0mu plus 1mu
\emergencystretch=3em
\setstretch{1.12}
\setlength{\parindent}{1.5em}
\setlength{\parskip}{0pt}

\title{Unrefusable Prompts and Assigned Credit in Consumer Software}
\author{Flyxion\\Independent Researcher}
\date{November 2026}

\begin{document}
\maketitle

\begin{abstract}
\noindent Two interface behaviors in widely used software are usually discussed separately. The first is the prompt that returns after it has been dismissed, exemplified by the Windows 11 backup and OneDrive reminders, which for years offered only a postponement as the negative answer. The second is the credit that a tool writes into a record without having been asked, exemplified by the co-author trailer that Claude Code adds to commits and, more weakly, by the default filename given to images downloaded from ChatGPT. This essay argues that in both cases a default determines which account of the interaction survives it, but that the two cases call for different remedies. The Windows prompt should let a person decline future reminders and keep that decision once it is made. The attribution should give the person a say before a byline is attached to their work, and should describe what the tool did instead of naming the tool as an author. The essay reviews what is publicly documented about each behavior, separates what the sources establish from what remains a user report, and notes that the same software identity can appear as a co-author in one interface and as a ranked developer in another through independent mechanisms. It does not establish that any vendor intended these outcomes, and it records that Microsoft has begun to add a permanent opt out to one of its prompts.
\end{abstract}

\vspace{2.5em}
\noindent\textbf{Keywords:} dark patterns, persistent prompts, default attribution, provenance, consent, authorship, software interfaces

\newpage

\section{Introduction}

A dark pattern is usually understood as an interface element that steers a person toward an action the operator prefers, such as a prominent button for acceptance beside a buried path for refusal. The examples considered here are less visible than a misleading button because they concern what the system keeps after the interaction ends. A prompt whose only negative answer is postponement gives a person who wants to decline indefinitely no way to say so. A trailer that names a product as a co-author of a commit has written a claim of authorship into a record that nobody was asked to approve. In both cases the interface settles a question that belongs to the person using it, and it settles that question through a default rather than through any exchange.

The two behaviors are treated together because a default decides, in each, what survives the interaction. The ``ask me later'' button writes an instruction to return, and that instruction outlasts the person's wish to be left alone. The co-author trailer writes a name into a record that travels with the work, and that name outlasts the person's wish to describe the assistance in their own terms. Beyond this the cases differ, and the difference matters because the remedies are distinct. The Windows case concerns a person who was offered postponement and wanted lasting refusal, so the interface should let the person state that decision and keep it once it is made. The attribution case concerns a person who may never have been asked whether they wanted the credit recorded at all, so the interface should give the person a say before it attaches a byline to their work.

A person who declines a backup reminder may want that decision to outlast the prompt. A person whose tool prepares a commit may want the resulting record to describe the assistance without assigning the tool a byline. In both cases the interface determines which account of the interaction will survive it. That decision matters especially when the record leaves the interface and becomes something other people rely on.

The argument proceeds from the documented behavior of each system to the two remedies. Care is taken throughout to distinguish three kinds of statement: what a vendor has documented, what independent reporters have observed on particular machines, and what individual users have reported. Claims of the last kind are marked as such.

\section{The Prompt That Returns}

Windows 11 has included a reminder, presented in several forms, that encourages the use of Windows Backup and, through it, OneDrive. Reporting in 2025 described a banner in the Settings application advising the person to back up the computer, and a commentator on the matter observed that the banner might be acceptable as a reminder provided there were a way to dismiss it and keep it from returning \cite{posey2025}. Users on Microsoft's own support forum asked how to stop the prompt from continually returning to the Settings home screen after every backup setting they could find had been disabled, and one of them asked whether the difficulty was a business policy intended to make most people give up \cite{techcommunity2025}. A separate trade report noted that users on Windows 11 Pro were urging a clearer opt out and that Microsoft had not said whether a group policy setting existed to suppress the messages \cite{windowsreport2025}.

The behavior has since changed in one respect. In 2026 Windows Latest reported that a full-screen backup prompt on stable Windows 11 25H2 builds now offers an option labelled ``Opt out of backup,'' and that the machines on which it was observed had no OneDrive sync client installed even though Windows was configured to back up files, folders, settings and credentials through Windows Backup \cite{elevenforum2026}. Secondary coverage described the previous situation plainly: the earlier reminder option settled nothing, because Windows could bring the prompt back later, leaving people repeatedly postponing something they had no intention of enabling \cite{digitaltrends2026}. A further report described the new button as the way to dismiss the prompt for good, and it did not claim that Microsoft had confirmed the rollout \cite{pcmag2026}. The same period saw a second report observing that a different Windows Backup reminder was flagged as required and could not be avoided, in contrast to another suggestion screen that could be switched off in the notification settings \cite{techradar2026}.

Three observations follow from these sources. The first is that the reappearance of the prompt was a property of how the negative answer was handled and not of any technical limit, because a button that records a refusal is among the simplest elements an interface can contain, and its later appearance in one prompt shows that the earlier omission could be reversed. The second is that the reports describe a return after a few days and give no fixed interval, so any precise figure for the period should be treated as a personal observation. The third is that the correction is partial. It applies to one prompt on one build family, and other suggestion surfaces remain governed by different rules.

What the sources support is a claim about what the interface allowed people to express. For years the only negative answer on offer was a postponement, so a person who wanted a lasting refusal had no way to give it from the prompt itself. The cost to that person was attention. Each return of the prompt required a decision already made, and the cumulative cost over the life of a machine is large relative to the single moment of effort that a persistent refusal would have required. The cost to the operator of recording the refusal is close to zero. An operator who measures conversions and does not measure irritation will drift toward the design that returns until the person gives in, and this explanation needs no assumption about intent, though it remains a hypothesis that the sources are consistent with and do not show.

\section{The Credit That Is Written}

A second behavior concerns what the system writes on the work itself, and it comes in a weaker and a stronger form.

The weaker form is the default filename given to images downloaded from ChatGPT. OpenAI documents that a person can upload an existing image and describe the changes to be made, and that ChatGPT Images generates the result \cite{openaihelp2026}. A downloaded file bearing the name of the product is accurate in one sense, since the product did produce the version that was downloaded. The difficulty is ambiguous provenance. The name identifies the tool that produced the downloaded file and says very little about who made the original, how much was changed, or whether the recipient is looking at a person's composition with a small edit or an image generated from a short description. A single default name covers all of these cases, and every later reader, recipient or platform inherits the flattened label. A third-party tool for removing metadata from such images notes that the default filename can reveal the origin of the file even after the embedded metadata has been stripped \cite{nanoglobals}, which shows that the name does travel with the work. The claim is limited to this, and it does not depend on how the model processes an upload internally.

The stronger form is the co-author trailer. Claude Code appends a line naming the model as a co-author to the commits it creates, and a pull request description receives a footer naming the tool \cite{classmethod2026}. A filename is a label that a reader may interpret. The trailer is a statement that the hosting platform interprets for them, because GitHub treats a co-author trailer as a designation of authorship, so the account named in it is regarded as one of the commit's authors and appears among the contributors of the repository \cite{classmethod2026}. The behavior was introduced as a default. A feature request in February 2026 stated that there was no way to disable it, and was closed within days \cite{issue22585}. The setting that came to control it was later renamed and its older form deprecated, and users have reported that the older setting was intermittently not respected \cite{issue4224,issue7543}. Another issue argues that the attribution should be opt in, on the ground that it adds information to commit messages that the developer did not author \cite{issue27083}.

The cost of reversing such a default is visible in an ordinary case. In the author's own cosmology repository, the intent was for the author record to identify the people who did the work and took responsibility for it. Achieving that required checking the Git history and GitHub's contributor display in order to remove the name of a commercial product, and the attention required to undo the credit was out of proportion to the ease with which it had been added. A credit that has entered a public record is expensive to withdraw, because the history, the contributor list and any copies made in the meantime must each be corrected.

A further consideration concerns who answers for a claim once it appears in the work. An earlier essay by the author distinguishes a reliable manner, which is a fact about outputs, from a conviction, which is a fact about what a party holds and defends, and it treats a consistent manner as weak evidence of conviction \cite{personbehind}. A chatbot can supply a persuasive defense of either position on request, and a human author can do the same and may have good reasons to. The difference lies in what follows publication. A person who publishes a claim under a byline can later be asked why it was made, what would lead to revising it, and what was done about an error that someone pointed out. When a coding tool drafts a change for a repository, a maintainer can ask it to generate an explanation, and the maintainer still has to decide whether that explanation holds up and whether the change belongs in the project. The byline is part of how a reader knows whom to ask, and a co-author line naming the tool points the reader toward a party that has not undertaken to answer for the claim.

\section{One Identity in Two Interfaces}

The trailer and the ranking are separate mechanisms, and the argument is stronger when they are kept apart. A follower-ranked table of accounts does not read commit trailers. It orders accounts by the number of people who follow them. What the two surfaces share is an identity: the same software account can appear as a co-author in the history of a repository and, independently, as a developer in a table sorted by followers. Each rendering is reasonable to the software that produces it. A software account on a hosting platform is an ordinary convenience, since something must post the results of a test run, open a pull request when a dependency changes, or publish a release, and in those contexts the account is visibly tied to an operation.

The difficulty arises when the identity moves between contexts in which different readers draw different inferences, and the history of the handle shows how much a current label can hide. Internet Archive captures of the profile at github.com/claude show a change in its public identity. On 12 May 2008 the handle was visible and active, with forks of two projects, commits to a small repository of its own, and four accounts followed \cite{wayback2008}. On 25 April 2019 it showed five public repositories, most of them forks of Ruby projects, four followers, and a photograph of clouds as its avatar \cite{wayback2019}. On 30 March 2023 it still showed what appears to be an individual's profile, with a personal display name, a personal website, five followers and six public repositories \cite{wayback2023}. On 23 May 2025 it displayed the name Claude, the product's logo as its avatar, a link to the Claude Code page, no public repositories and 81 followers \cite{wayback2025}. A handle that once displayed an individual's profile now displays a commercial product, and the ranking presents the current identity as though it were simply another developer. The captures establish a change in the profile's public identity between 30 March 2023 and 23 May 2025. They do not establish how control of the handle passed, and no claim about that is made here.

The ranking then places the two identities side by side. At the time of observation, a follower-ranked table of accounts on the platform placed the Claude account directly beneath the account of the person who initiated Linux and created Git \cite{leaderboard}. The profile that had 81 followers in the May 2025 capture appeared in the November 2026 ranking directly below Torvalds, in rows that shared the same fields, the same sort of avatar and the same numerical measure of public standing. The table gives both rows equal precision and says nothing about how either audience accumulated. Two days after the ranking cache was created, the two profile pages displayed 239 thousand followers for the Claude account and 332 thousand for the Torvalds account \cite{profiles}, so the first account had gained about three thousand followers in the interval while the second showed no change at the precision the page displays. Only two readings exist and both are rounded, so they indicate direction and not a rate.

Two dates for the account should be kept apart. The platform's API reports a creation date of 7 May 2009 for the current record of the Claude account and 3 September 2011 for the Torvalds account \cite{apidata}. The 2008 capture shows that the handle was visible and active before the 2009 date, and the API date applies to the current account record. The reason the API date is later than the earliest archived activity is not documented in the sources consulted, and neither date is used here to settle the other. The handle was in use before the Torvalds account was created and before Anthropic was founded in 2021 \cite{anthropicwiki}.

No party needed to plan this outcome. The platform supplies the account and its attribution machinery, the vendor supplies the name, and a ranking site sorts by a count. The finished page makes a claim that none of them had to formulate, which is that the product is to be understood in the same terms as the people listed around it. The co-author trailers do not cause the ranking, but they contribute to the impression that the account is a participant in the work and not a tool used in it.

Behind that single name is a large body of collective work. People wrote libraries, answered questions, maintained documentation, posted examples and explained obscure failures in issue threads, often so that the next person would have an easier time. A coding model draws much of its usefulness from this accumulated practice. Where an output reproduces identifiable code there is a concrete question about the source and the acknowledgment owed to its author. In most cases the debt is harder to trace, because patterns learned across many projects do not yield a clean account of who contributed what to a particular answer. The difficulty of distributing credit across all those contributors becomes, in the interface, the ease of giving it to one brand. The point does not depend on the legal status of training. It depends on the observation that the interface offers a place for exactly one name.

\section{Two Defaults, Two Remedies}

The Windows prompt and the attribution trailer are both defaults that decide what survives the interaction, but they fail in different ways, and a single remedy would fit neither well.

The prompt fails to let the person decline. A person who has postponed it has given an answer, and Windows could act on that answer exactly as designed, but a person who wanted to decline indefinitely had no way to say so. The interface should offer that answer and keep it once it is given. A dismissal should therefore offer at least three answers that are stored: ask me later, never ask me about this, and never ask me about anything in this category. The second and third require one control each and no new infrastructure. The reminder interval for the first should be disclosed, and the settings that control a category of suggestions should be reachable from the prompt itself and not from a different part of the system found only through search. The fault in the existing design is an omission in what the interface let people say, and a permanent opt out on one prompt, as in the recent change, repairs it for that prompt only.

The attribution fails to ask. A person whose tool prepares a commit has made no statement about credit, and the interface fills the silence with the tool's own name. The remedy is therefore to give the person a say before a byline is attached, and to change what is written when the person says nothing. A commit could record the tool and the operation, for example that a test suite was run, that a function was drafted, or that a file was reformatted, and leave the author field to the people who reviewed and accepted the change. Identifiable sources would keep their acknowledgments. Where a person wishes to credit the tool as a co-author, the trailer remains available, and the argument is for symmetry and not for removing the feature. A downloaded image could likewise carry a neutral name by default, with the choice of a name that reflects its relationship to any upload made available before the file is written. A reader would then have something to go on beyond the fact that a particular vendor's software was present.

Disclosure of the use of software is useful and is not in dispute. If a person uses a program to check a proof, format a bibliography or revise a function, the program's role can be stated, and a reader can decide whether it matters. A co-author line says much less. It does not tell the reader which parts were generated, which suggestions were discarded, what was verified, or how much work went into making the result usable. What it does very efficiently is place the product's name next to the person's. The ranking surface admits a related repair, since a table that sorts accounts by followers could separate accounts that act as tools from accounts that represent people, and the platform already holds the information needed to do so.

\section{Limits of the Argument}

Several limits apply. The essay does not establish intent on the part of any vendor. The reports on the Windows prompts come from trade coverage and user forums, and Microsoft has not been shown to have confirmed the recent change. The claim about image filenames is restricted to what the vendor documents and what the file carries, and the essay makes no claim about how images are processed internally or about the naming of files that are returned unchanged. The behavior of the ranking depends on the date of observation and on the site that produced it. Some people want the trailer, the filename or the reminder, and a design that stores a refusal should store an acceptance equally. Finally, the essay does not measure how often any of the behaviors occurs, and the scale of the problem is asserted from the structure of the interfaces and not from counts.

\section{Conclusion}

A prompt that returns and a credit that is written are both ways of deciding, without an exchange, which account of an interaction will last. A backup prompt should let a person decline future reminders and keep that decision once it is made. A coding tool should ask before adding a byline and describe its contribution in terms a reader can use. The Windows case concerns what a person is allowed to say to the system. The attribution case concerns what the system writes into a public record on the person's behalf. The partial correction to one Windows prompt shows that the first remedy is within reach, and the existence of a configurable attribution setting shows that the second is as well. What remains is to make the better answer the default and not the exception.

\clearpage
\begin{thebibliography}{99}
\raggedright

\bibitem{posey2025} Posey, B. Microsoft's Latest Backup Prompt May Cause Unforeseen Problems. \emph{Redmond Magazine}, June 2025. \url{https://redmondmag.com/articles/2025/06/06/microsoft-latest-backup-prompt-may-cause-unforeseen-problems.aspx}

\bibitem{techcommunity2025} Microsoft Tech Community Forum Discussion. Action Advised: Back Up Your PC. \url{https://techcommunity.microsoft.com/discussions/windows11/action-advised---back-up-your-pc/4397852}

\bibitem{windowsreport2025} Windows Report. Windows 11 Backup Prompt Can Disrupt Your Workflow. \url{https://windowsreport.com/windows-11-backup-prompt-can-disrupt-your-workflow/}

\bibitem{elevenforum2026} Windows 11 Will Now Let You Opt Out of Backup to OneDrive. Eleven Forum thread quoting Windows Latest, 2026. \url{https://www.elevenforum.com/t/windows-11-will-now-let-you-opt-out-of-backup-to-onedrive.49119/}

\bibitem{digitaltrends2026} Digital Trends. Windows 11's Needy OneDrive Nagging Is About to Go Away. 2026. \url{https://www.digitaltrends.com/computing/windows-11s-needy-onedrive-nagging-is-about-to-go-away/}

\bibitem{pcmag2026} Microsoft Finally Lets You Dismiss the OneDrive Prompts for Good. Republished from PC Magazine, APH Networks, 2026. \url{https://aphnetworks.com/news/32100-microsoft-finally-lets-you-dismiss-onedrive-prompts-good}

\bibitem{techradar2026} TechRadar. Microsoft Finally Fixes One Frustrating Windows 11 Pop-Up, Just as a More Irritating One Arrives to Take Its Place. 2026. \url{https://www.techradar.com/computing/windows/microsoft-finally-fixes-one-frustrating-windows-11-pop-up-just-as-a-more-irritating-one-arrives-to-take-its-place}

\bibitem{openaihelp2026} OpenAI. ChatGPT Image Library. OpenAI Help Center. \url{https://help.openai.com/en/articles/11084440-chatgpt-image-library}

\bibitem{nanoglobals} Remove Metadata from ChatGPT Images. NanoGlobals tool page. \url{https://nanoglobals.com/tools/metadata-remover/chatgpt/}

\bibitem{classmethod2026} I Tried Hiding the Commit and Pull Request Attribution for Claude Code. Classmethod Developers Blog. \url{https://dev.classmethod.jp/en/articles/claude-code-attribution-setting/}

\bibitem{issue22585} Feature Request: Optional Co-Authored-By Attribution. GitHub issue 22585, anthropics/claude-code, February 2026. \url{https://github.com/anthropics/claude-code/issues/22585}

\bibitem{issue4224} Settings for includeCoAuthoredBy and gitAttribution Intermittently Not Respected During Commits. GitHub issue 4224, anthropics/claude-code. \url{https://github.com/anthropics/claude-code/issues/4224}

\bibitem{issue7543} includeCoAuthoredBy False Setting Not Respected. GitHub issue 7543, anthropics/claude-code, September 2025. \url{https://github.com/anthropics/claude-code/issues/7543}

\bibitem{issue27083} Co-Authored-By Attribution Should Be Opt-In, Not Opt-Out. GitHub issue 27083, anthropics/claude-code. \url{https://github.com/anthropics/claude-code/issues/27083}

\bibitem{leaderboard} GitHub Rank. Follower-ranked table of GitHub accounts, cache created 15 November 2026, data from the GitHub API v3. \url{https://wangchujiang.com/github-rank/}

\bibitem{profiles} GitHub profile pages for the accounts claude and torvalds, viewed by the author on 17 November 2026 at about 9:51 and 9:52 Atlantic time. \url{https://github.com/claude} and \url{https://github.com/torvalds}

\bibitem{apidata} GitHub REST API user records for claude (created 2009-05-07, public repositories 0) and the torvalds profile activity overview (joined 3 September 2011), retrieved by the author on 17 November 2026. \url{https://api.github.com/users/claude}

\bibitem{anthropicwiki} Anthropic. Wikipedia. \url{https://en.wikipedia.org/wiki/Anthropic}

\bibitem{wayback2019} Internet Archive Wayback Machine. Capture of the GitHub profile github.com/claude, 25 April 2019. \url{https://web.archive.org/web/20190425022604/https://github.com/claude}

\bibitem{wayback2008} Internet Archive Wayback Machine. Capture of the GitHub profile github.com/claude, 12 May 2008. \url{https://web.archive.org/web/20080512160629/https://github.com/claude}

\bibitem{wayback2023} Internet Archive Wayback Machine. Capture of the GitHub profile github.com/claude, 30 March 2023. \url{https://web.archive.org/web/20230330132637/https://github.com/claude}

\bibitem{wayback2025} Internet Archive Wayback Machine. Capture of the GitHub profile github.com/claude, 23 May 2025. \url{https://web.archive.org/web/20250523110530/https://github.com/claude}

\bibitem{personbehind} Flyxion. The Person Behind the Interface: Authorship, Identity, and the Attribution of Mind to AI. October 2026. \url{https://standardgalactic.github.io/antivenom/mind/person-behind-the-interface.pdf}

\end{thebibliography}

\end{document}
